\documentclass[10pt,a4paper]{article}

\usepackage[utf8]{inputenc}
\usepackage[cyr]{aeguill}
\usepackage[francais]{babel}

% Les packages
%\usepackage[french]{babel}
\usepackage{fancybox}
\usepackage{graphics}
\usepackage{color}
\usepackage{latexsym}
\usepackage{amsmath,amsthm,amssymb}
\usepackage[plain]{fullpage} 
\usepackage{verbatim}
\usepackage{times,epsfig}
\usepackage{listings}
\usepackage{multirow}
\usepackage{amsmath,amsthm,amssymb,amscd,txfonts,bbm} % RAJOUT ICI !!!
\usepackage{graphics,epsfig} % RAJOUT ICI !!!
%\usepackage{algorithmic}
\usepackage[final]{pdfpages} 
\usepackage{fancyvrb}
\usepackage{enumerate}
\usepackage{listingsutf8}
\usepackage{comment}
\usepackage{minted}

% small et verbatim: to be coherent with \file
\newenvironment{smallverbatim}%
{\endgraf\small\verbatim}{\endverbatim}

\newenvironment{qv}
{\quote\Verbatim}
{\endquote\endVerbatim}

\usepackage{lastpage}

\newcommand{\terme}[1]{\texttt{#1}}
\newcommand{\termSet}[1]{\{\texttt{#1}\}}
\newcommand{\guill}[1]{<<\,#1\,>>}
% Macro pour l'inclusion de figure  
\newcommand{\fig}[2]{%
  \begin{figure*}[hbt]
   \begin{center}
      \includegraphics[width=6cm]{#1}
  \caption{\label{#1}  #2}
  \end{center}
 \end{figure*}
}
\newcommand{\figTex}[2]{%
  \begin{figure*}[hbt]
 \centerline{\input{#1.pstex_t}}
  \caption{\label{#1}  #2}
 \end{figure*}
}

% Macro commentaires
\definecolor{turquoise}{rgb}{0.20 ,0.60, 0.60}
\definecolor{vert}{rgb}{0.40,  0.60 ,0.00}
\definecolor{violet}{rgb}{0.80 , 0.00 , 0.80}
\definecolor{bleu}{rgb}{0.20,  0.20,  0.80}
\definecolor{rose}{rgb}{1.00,  0.20,  0.40}

\newcommand{\cle}[1]{\textbf{#1}}
\def\fk#1{\textit{#1}}
\def\tble#1{\textit{#1}}

\lstset{basicstyle=\footnotesize,showstringspaces=false,language=XML,breaklines=false,columns=fullflexible,frame=shadowbox, captionpos=b}

\newenvironment{solution}[0]{
~\\ \color{red} \textbf{Solution : }\color{black}}{%
}
\excludecomment{solution}


\parindent=0pt           
                             
\usepackage{url}

\newtheorem{Exercice}{Definition}


\begin{document}

\title{Sujet d'examen UE NFE204}
\author{Année universitaire 2023-2024}
\date{Examen première session: 22 juin 2024}

\maketitle

\begin{center}
{\bf \Large Consignes}\\
\begin{tabular}{| p{15cm} |}
\hline
Dur\'ee de l'examen : 3h00; Calculatrice autoris\'ee\,; Documents autorisés: 5 pages recto-verso de note personnelles.\\

Les exercices sont indépendants les uns des autres. Merci de soigner votre écriture.\\
Ce sujet contient \pageref{LastPage} pages.\\

\textit{Important}\,: nous n'évaluons pas les réponses par la syntaxe. Ne vous
inquiétez pas d'utiliser un point-virgule à la place d'un point d'exclamation ou
autres détails mineurs
\\
\hline
\end{tabular}
\end{center}

\section*{Première partie : documents structurés (6 pts)}

Nous approchons des jeux olympiques de Paris 2024, événement majeur,
mais source probable de désordres importants. Les organisateurs 
décident de mettre en place un système de co-voiturage à destination
des sites majeurs, afin d'éviter les encombrements.
 
On pense au départ d'une base relationnelle sera suffisante, et on
adopte le modèle suivant (très simplifié bien sûr). Les clés
primaires sont en \textbf{gras}, les clés étrangères en 
\textit{italiques}.

\begin{itemize}
  \item Personne   (\textbf{id}, nom, domicile)
  \item Site (\textbf{id}, nom)
  \item Trajet (\textbf{id}, \textit{idConducteur}, \textit{idSite}, date)
  \item Passager (\textbf{\textit{idTrajet}, \textit{idPassager}})
\end{itemize}

\vspace*{0.2cm}

Voici également un tout petit échantillon de la base.

\vspace*{0.2cm}
\begin{center}
\begin{minipage}[b]{4cm}
\begin{tabular}{l|l|l}
\cle{id} & nom & domicile\\
\hline
1 & Albert & Melun\\
2 & Aline & Versailles\\
3 & Ktie & Créteil\\
\hline
\end{tabular}

\centerline{La table \tble{Personne}}
\end{minipage} \hspace*{0.5cm} \
\begin{minipage}[b]{3cm}
\begin{tabular}{l|l}
\cle{id} &  nom  \\
\hline
100 & Stade de France\\
101 & Parc de Versailles \\
102 & Grand Palais \\
\hline
\end{tabular}

\centerline{La table \tble{Site}}
\end{minipage} 

\vspace*{0.5cm}
\begin{minipage}[b]{6cm}
\begin{tabular}{l|l|l|l|l}
\cle{id} & idConducteur & idSite & date \\
\hline
A & 1 & 100 & 30/07 \\
B & 3 & 100 & 4/08 \\
C & 1 & 101 & 31/07 \\
\hline
\end{tabular}

\centerline{La table \tble{Trajet}}
\end{minipage} \hspace*{0.5cm} \
\begin{minipage}[b]{3cm}
\begin{tabular}{l|l|l}
\cle{id} &  idPassager  \\
\hline
A & 2 \\
A & 3 \\
B & 1 \\
C & 2 \\
\hline
\end{tabular}

\centerline{La table \tble{Passager}}
\end{minipage} 
\end{center}

\vspace*{0.2cm}

On s'aperçoit rapidement que cette représentation impose 
beaucoup de jointures et ne passe pas à l'échelle. 
On décide de remplacer la base relationnelle par une base Cassandra. 



\begin{enumerate}
  \item Voici une première table Cassandra pour représenter les trajets.
  
\begin{minted}{sql}
create table Trajet1 (idTrajet text, 
                    date date, 
                    site frozen<Site>, 
                    conducteur frozen<Personne>, 
                    passagers set< frozen<Personne>>,
                 primary key idTrajet );
\end{minted}

 Donnez l'exemple d'un  document JSON correspondant à la structure,
           de la table \texttt{Trajet1},
           avec les informations du trajet A de la base relationnelle ci-dessus.
           
  \item Voici la proposition d'une autre modélisation, avec une
    table \texttt{Trajet2} et une table \texttt{Passager} dont la clé est
      \emph{composite}
    
\begin{minted}{sql}
create table Trajet2 (idTrajet text, 
                    date date, 
                    site frozen<Site>, 
                    conducteur frozen<Personne>
                 primary key (idTrajet);

create table Passager (idTrajet text, 
                     idPassager text,
                    passager frozen<Personne>,
                 primary key (idTrajet, idPassager );
\end{minted}

Rappelons ce qu'est une clé composite en Cassandra, d'après la documentation\,: 
\textit{A compound primary key consists of the partition key and the clustering key. The partition key determines which node stores the data. Rows 
for a partition key are stored in order based on the clustering key.} Quelles affirmations
sont vraies\,?
\begin{enumerate}
  \item Les lignes de \texttt{Trajet2} et les lignes de \texttt{Passager} sont sur le même serveur
      \item Les passagers d'un même trajet sont tous sur un seul serveur 
       \item Les passagers stockés sur un serveur sont triés sur leur identifiant
       \item Les passagers d'un même trajet sont stockés contiguement
\end{enumerate}  

  \item Est-il possible de convertir la table \texttt{Trajet1} 
  vers la table \texttt{Passager} avec un traitement
        Map Reduce? Expliquez comment, en indiquant ce que font les fonctions de Map et de Reduce.
\end{enumerate}

\section*{Deuxième partie : recherche d'information (6 pts)}

L'application enregistre des commentaires déposés par les clients sur les conducteurs, et réciproquement. 
Voici un échantillon.

\vspace*{0.3cm}

\begin{tabular}{lp{15cm}}
     $c_1$& Conduite dangereuse, et conducteur bavard. \`A éviter.\\
     $c_2$& Pour la conduite ça ne va pas\,: le conducteur (sympa par ailleurs)
         réussit à être à la fois lent et dangereux\\
     $c_3$& Trop dangereux\,! Je veux bien payer pour un co-voiturage, mais
          pas avec un conducteur dangereux.\\
     $c_4$& Sympa, mais le conducteur est très bavard et lent.\\
     $c_5$& Il est vraiment bavard de chez bavard. Heureusement, sympa et rien
         de dangereux dans sa conduite.\\ 
\end{tabular}
             
On va se limiter au vocabulaire (\texttt{dangereux, bavard, lent, sympa}).
(NB: on suppose
     une phase préalable de normalisation qui élimine les pluriels, majuscules, trouve
     que ``sympa'' et ``sympathique'', ``dangereux'' et ``dangereuse'',  sont les mêmes mots, etc.)

\begin{enumerate}
   \item Donnez une matrice d'incidence contenant les tf, 
      et un tableau donnant les idf (sans le $\log$), pour les  termes précédents. Placez
      les termes en ligne, les commentaires en colonne.

\begin{solution}
\begin{center}
\begin{tabular}{l|c|c|c|c|c}
            & \textbf{c1} & \textbf{c2} & \textbf{c3} &\textbf{c4}&\textbf{c5} \\ \hline  
   dangereux  5/4   & 1           & 1           & 2           & 0         & 1\\
   bavard 5/3  & 1           & 0           & 0           & 1         & 2\\
   lent 5/2  & 0           & 1           & 0           & 1         & 0\\
   sympa 5/3  & 0           & 1           & 0           & 1         & 1\\\hline
   \end{tabular}
\end{center}
\end{solution}

  \item Donnez les normes des vecteurs représentant ces commentaires.
\begin{solution}
Les normes
\begin{itemize}
  \item $||c_1|| = \sqrt{1+ 1}$; $||c_2|| = \sqrt{3}$; $||c_3|| = \sqrt{4}$; 
    $||c_4|| = \sqrt{3}$, $||c_5|| = \sqrt{6}$  
\end{itemize}
\end{solution}
  
  \item Donner les résultats classés par similarité cosinus basée sur les tf 
  (on ignore l'idf) pour les requêtes suivantes. Expliquez brièvement la raison
  du classement pour le premier document.
  
  \begin{itemize}
    \item bavard;
    \item lent et sympa;
    \item dangereux et bavard.
  \end{itemize}
\begin{solution}
Les cosinus (requête non normalisée).

\begin{itemize}
  \item bavard\,: c1 : $\frac{1}{\sqrt{2}}$;  
       c4 : $\frac{1}{\sqrt{3}}$; 
       c5 : $\frac{2}{\sqrt{6}}$; 
  
  \item lent et sympa\,: 
  
    c2 : $\frac{1+1}{\sqrt{3}}$; 
    c4 :  $\frac{1+1}{\sqrt{3}}$  
    c5 : $\frac{1}{\sqrt{6}}$; 

  \item dangereux et bavard
  
    c1 : $\frac{1+1}{\sqrt{2}}$; 
    c2 : $\frac{1}{\sqrt{3}}$;
    c3 :  $\frac{2}{\sqrt{3}}$  
    c4 :  $\frac{1}{\sqrt{3}}$  
    c5 : $\frac{1+2}{\sqrt{6}}$; 
    
 \end{itemize}

\end{solution}  

\begin{comment}
   \item Pour la dernière requête, ``dangereux et bavard'', on constate que deux
     documents sont à égalité. Est-ce que cela change 
   si on prend en compte l'idf, comment et pourquoi\,?
   
\begin{solution}
bavard est un peu moins présent que dangereux dans la collection, ce qui amène à classer
c4 avant c2.
\end{solution}
\end{comment}
   \item Pour la requête ``lent et sympa'', et les documents c2 et c4,
   qu'est-ce qui change si on ne met pas de restriction 
   sur le vocabulaire (tous les mots sont indexés)?
   
\begin{solution}
Dans ce cas le plus petit document (c4) l'emporte car sa norme est
plus petite. Intuititvement: il parle plus \emph{spécifiquement} de
lent et de sympa que c2.
\end{solution}   
    
%\item On voudrait pouvoir mesurer les co-occurrences de termes, autrement dit
%    produire un indicateur représentant la probabilité que deux termes
%    apparaissent ensembles dans un même commentaire. Proposez une formule estimant
%    cette probablilité et s'appuyant
%      sur les informations présentes dans la matrice d'incidence.

\end{enumerate}


\section*{Troisième partie : MapReduce (3 pts)}

Nous recevons un flux de commentaires dans
un fichier \texttt{trajets.csv}. Il contient l'identifiant
du conducteur, du client, et le commentaire. Voici le début.

\begin{smallverbatim}
1, 101, "Voiture bruyante"
1, 103, "Une vraie épave"
2, 101, "La voiture est propre, plus que le conducteur..."
..
\end{smallverbatim}

Utilisant Pig latin, on charge cette collection avec la commande suivante\,:
\begin{smallverbatim}
trajets = LOAD 'trajets.csv' as (idConducteur, idClient, commentaire);
\end{smallverbatim}

Et on exécute le script Pig suivant\,:

\begin{smallverbatim}
coll1 = filter trajets by contains(commentaire, "voiture")
coll2 = group coll1 by idConducteur;
coll3  = foreach coll2 generate group, COUNT(coll1.commentaire);
\end{smallverbatim} 

\begin{enumerate}
 \item Expliquez le résultat produit par ce script.
 \item Comment ce script peut-il être traduit en MapReduce? Donnez la fonction
    de Map et la fonction de Reduce, sous la forme de votre choix.
\item Quelle serait la requête SQL correspondant à ce script Pig\,?

\end{enumerate}

\section*{Quatrième partie : systèmes distribués (3 pts)}

\begin{enumerate}
 \item Sachant que Cassandra organise la distribution des données 
 selon la technique du hachage cohérent, expliquez l'algorithme qui
 place un document de la table \texttt{Passager} sur un serveur.
  
  \item Parmi les requêtes ci-dessous, lesquelles peuvent être 
  routées vers un seul serveur\,?
  
  \begin{enumerate}
  \item Les trajets vers Versailles  
   \item Les passagers du trajet  $t12$ (c'est son identifiant)
    \item Les sites visités par le passager $p988$ (c'est son identifiant) 
\end{enumerate}  


  \item À l'instant $t$, les valeurs de hachage de deux trajets
  $T_1$ et $T_2$ sont identiques.  Peuvent-ils être placés sur deux
  serveurs différents à un instant $t+\Delta$, sachant qu'entretemps 
  plusieurs nœuds ont été ajoutés au système Cassandra\,?
\end{enumerate}


\section*{Cinquième partie : question de cours (2 pts)}


\begin{enumerate}
 \item Comment expliquer l'expression ``index inversé'' dans  
     la terminologie des moteurs de recherche\,?
  
 \item Rappeler le principe de \textit{data locality}
\end{enumerate}

\end{document}
